\begin{abstract}
The Simplified General Perturbations-4 (SGP4) propagator is the industry standard for satellite tracking, space traffic management (STM), and conjunction assessment. However, it suffers from significant accuracy degradation during geomagnetic storms due to unmodeled upper atmospheric expansion and dynamic drag coefficients. In this paper, we present a novel ''Gated Physics-Informed Neural Network'' (G-PINN) designed to correct SGP4 residual errors in Low Earth Orbit (LEO) dynamically. Unlike traditional multi-layer perceptrons (MLPs) or recurrent architectures that often introduce catastrophic ''phantom drift'' at initialization by predicting non-zero errors at the Two-Line Element (TLE) epoch, our architecture introduces a mathematically differentiable ''Physics Gate'' ($\tanh(\lambda t)$). This gate enforces a strict hard constraint, ensuring zero initial error while allowing the network to continuously apply its learned non-linear drag corrections over propagation time. Evaluated on a massive, highly curated dataset of historical Starlink trajectories (2022--2025) augmented with Celestrak space weather indices ($K_p$ and F10.7 solar flux), the G-PINN demonstrates a $12.6\%$ reduction in maximum position error during severe solar storms, reducing mean storm-induced spatial uncertainties from 11.18 km to 9.77 km. During quiet space weather periods, the network achieves absolute stability, doing ''no harm'' and maintaining SGP4 baseline integrity. Furthermore, we report a scientific discovery facilitated by the model's internal correction structure: the identification of a systematic $5.5$ m/s radial velocity bias in the SGP4 target handling, effectively allowing the PINN to reverse-engineer and flag degraded ballistic $B^*$ terms in public TLEs. Finally, we implement a robust 3D Mahalanobis-distance-based Monte Carlo conjunction assessor and a WebGL-based high-fidelity orbital visualization engine to transition these AI corrections into actionable, real-time risk classification systems.
\end{abstract}
